\clearpage\markboth{}{}
\adjustmtc
\chapter*{Conclusion}
\label{cha:conclusion}
\addstarredchapter{Conclusion}



\section*{Contributions}
\newcommand\sectionname{Contributions}
\markboth{Conclusion}{\sectionname}


La réévaluation de modèles paramétriques est une fonctionnalité incontournable
des modeleurs interactifs. %
Ceux-ci permettent en effet d'éviter de refaire le travail fastidieux de
conception d'un modèle pour en générer une simple variation. %
Cela s'avère particulièrement avantageux lorsque le processus de conception ne
change pas de manière substantielle~; par exemple, avec la modification des
paramètres géométriques des opérations qui n'implique aucun ajout, suppression
ou déplacement d'opération. %

Afin d'être considéré fonctionnel, un mécanisme de réévaluation doit respecter
un certain nombre de contraintes. %
Le mécanisme de nomination doit permettre l'identification de manière unique et
non ambiguë des entités topologiques. %
L'objet issu du processus de réévaluation doit respecter les intentions de
modélisation de l'utilisateur, c'est-à-dire que le programme ne doit pas prendre
d'initiative à la place de l'utilisateur. %
L'objet doit être topologiquement et géométriquement cohérent. %

Dans les solutions existantes, nous avons identifié plusieurs faiblesses liées à
ces différentes contraintes. %
Premièrement, si les systèmes de nomination permettent déjà d'identifier des
entités topologiques avec une certaine fidélité, ceux-ci ne sont pas
génériques. %
Leurs schémas de nomination sont adaptés à une dimension précise de
modélisation. %
Les étendre à d'autres dimensions nécessite alors des efforts importants ainsi
que la définition de nouvelles procédures d'appariement. %
Deuxièmement, le recours à la géométrie de certaines méthodes introduit des
ambiguïtés lors de l'appariement. %
Ces ambiguïtés entraînent elles-mêmes des erreurs d'appariement et, par
conséquent, le non-respect des intentions de modélisation exprimées par un
utilisateur. %
Enfin, ces mêmes erreurs d'appariement peuvent produire des objets qui ne sont
pas cohérents du point de vue topologique. %
Il est donc nécessaire de contrôler et, le cas échéant, de réparer l'objet
produit. %



\paragraph*{Détection automatique des évènements}



Notre première contribution concerne la détection des évènements (création,
suppression, scission, fusion, non-changement et modification). %
Nous avons formalisé un ensemble de critères permettant de caractériser chacun
des évènements aussi bien dans les objets que dans les règles de transformation
de graphe. %

Nous avons ensuite étendu ce formalisme au langage de règle Jerboa qui permet de
représenter des opérations de modélisation géométrique. %
De plus, l'analyse étant effectuée uniquement à partir des règles Jerboa, les
évènements sont déterminés indépendamment des objets sur lesquels les règles
sont appliquées. %
Nous décrivons aussi les limites des détections liées à ce formalisme. %
Dans de tels cas, une simple vérification locale, dans l'objet, permet de
confirmer ou non la détection d'un évènement. %

Avec ce formalisme, nous proposons une nouvelle méthode permettant de suivre
l'évolution des entités topologiques dans un modèle, exclusivement grâce à une
analyse syntaxique des règles Jerboa. %
Celle-ci peut être opérée statiquement en pré-calcul au premier lancement du
modeleur, ce qui s'avère particulièrement efficace en temps de calcul. %
Ainsi, cette méthode s'affranchit de la nécessité de parcourir en parallèle
deux objets pour les comparer. %
De même, il devient inutile de spécifier explicitement quels évènements sont
attendus suite à l'application d'une opération. %

Enfin, cette méthode de détection est générique. %
Celle-ci étant basée sur l'analyse des règles, elle n'est pas conçue pour une
dimension spécifique. % %%MAXIME: revoir le découpage de cette phrase
Elle permet alors la détection des évènements quelle que soit la
dimension de modélisation du modeleur ou des orbites manipulées par l'opération. %



\paragraph*{Réévaluation dans un système de modélisation à base de règles}



À l'aide des formalismes des cartes généralisée et des règles Jerboa enrichis de
notre formalisation des évènements, nous avons implanté un mécanisme de
nomination persistante en deux couches. %
Ce mécanisme consiste, dans un premier temps, à construire puis maintenir un
historique d'opérations et de n\oe uds de règles pour chacun des brins d'une
carte généralisée. %
Grâce à cela, tout brin peut être identifié de manière unique et non ambiguë. %
Ainsi, à chaque fois qu'une règle est appliquée, l'historique du brin
sélectionné à l'instance courante de modélisation est enregistré dans une
spécification paramétrique. %
Cet historique ainsi enregistré est immuable et devient alors un nom persistant de
brin. %

Dans un deuxième temps, nous avons implanté une procédure qui, pour chaque nom
persistant enregistré au cours de la modélisation, permet la reconstitution de
l'historique des évolutions d'une orbite incidente au brin désigné. %
Un tel historique est représenté par un DAG dit d'évaluation nous permettant de
traiter l'information qu'il contient. %
Dans cette procédure, nous utilisons la formalisation d'évènements pour détecter
et enregistrer dans le DAG chacun des évènements impactant une orbite au cours
de la modélisation. %
De plus, cet historique est enrichi de la notion d'origine qui nous permet alors
d'identifier sans ambiguïté une orbite lors de la réévaluation sur la base de son
historique de construction. %
Ces origines sont elles-mêmes tracées jusqu'à représenter leurs ancêtres les plus
anciens. %

Pour la réévaluation à proprement parler, nous avons implanté une
procédure d'appariement qui crée de nouveaux DAG de réévaluation en reproduisant
d'abord les DAG d'évaluation puis en les mettant à jour. %
De cette manière, nous pouvons conserver à la fois les DAG d'évaluation et de
réévaluation. %
Les DAG de réévaluation permettent la mise en \oe uvre de procédures
d'appariement au fur et à mesure de leur construction. %
Ainsi, une fois le DAG est complété, celui-ci désigne les brins à utiliser comme
paramètres topologiques pour l'application suivante. %

Enfin, nous proposons de prendre en compte l'ajout, la suppression et le
déplacement d'opérations. %
Ce type de modification au niveau de l'historique de construction a un impact
plus important sur la réévaluation et les mécanismes
de nomination persistante que la simple édition des paramètres géométriques,
qui s'avère restrictive. %
En complément, nous proposons des stratégies paramétrables pour adapter la réévaluation. %
Par exemple, lorsqu'une orbite est scindée par une application ajoutée dans la
spécification réévaluée, il est possible de choisir les différents paramètres
topologiques sur lesquels poursuivre la réévaluation. %
Au contraire, lorsque deux orbites sont fusionnées par la suppression d'une
sous-orbite, nous proposons de choisir de préserver l'une ou
l'autre ou les deux historiques de ces orbites dans celle fusionnée. %
Actuellement, seuls l'ajout et la suppression d'opération sont supportés par le
mécanisme de réévaluation. % %%MAXIME: tout relire pour voir si ça ne casse pas la lecture



\paragraph*{Extension aux scripts}



Les règles Jerboa ne permettent pas de représenter des opérations complexes,
c'est-à-dire faisant appels à plusieurs orbites non isomorphes. %
Ces opérations complexes sont réalisées par les scripts Jerboa qui composent
plusieurs règles. %
Nous proposons une extension de notre mécanisme de réévaluation afin de
supporter les scripts Jerboa et donc ces opérations complexes. %
Dans un premier temps, nous définissons les extensions nécessaires à nos
structures de données pour trois des structures de contrôles usuelles~: la
séquence, l'itération de type \texttt{foreach} et l'alternative. %

Les scripts représentent, d'une certaine façon, des spécifications paramétriques
à part entière. %
Dans notre démarche, nous avons fait le choix de la représentation en «~boîte
ouverte~». %
En effet, contrairement à la «~boîte fermée~» qui permet seulement de gérer un
script comme une règle, la boîte ouverte nous permet d'accéder au contenu des
scripts, à savoir les opérations, les paramètres topologiques et géométriques
ainsi que les structures de contrôle utilisées. %

Nous proposons ainsi, sur la base de la représentation boîte ouverte, des
extensions de structures de spécification paramétrique et de DAG d'évaluation et
de réévaluation. %
% \xavier{[Remarque : pourquoi parler de boîte fermée ? Tout ton discours précédent et suivant repose sur la boîtes ouverte. Si tu veux simplement rappeler la différence entre les deux, fais-le en début de paragraphe, puis explique le choix de la boîte ouverte et ne parle plus de la boîte fermée.]}%
Dans les spécifications paramétriques, l'extension doit nous permettre
d'enregistrer des paramètres topologiques à la fois pour les scripts et pour les
règles qu'ils appellent. %
Il est également nécessaire d'enregistrer les structures de contrôles employées
afin de représenter fidèlement le processus constructif d'un objet. %
Dans les DAG, nous réutilisons les informations de scripts ajoutées dans la
spécification paramétriques. %
Ainsi, les opérations appartenant à un script sont englobées dans celui-ci. %
De plus, nous proposons de conditionner la réévaluation d'une orbite en fonction
des paramètres topologiques des scripts qui impactent son évolution ainsi qu'à
leurs structures de contrôles. %
De cette manière, une structure itérative est conditionnée par l'existence de
l'orbite itérée ainsi que celle de ses instances de sous-orbites d'itération. %

Enfin, nous avons implanté une nouvelle procédure d'appariement d'orbite dans le
cadre de l'alternative. %
Étant donné que deux règles peuvent être très différentes, nous proposons une
procédure capable d'apparier des orbites initialement évaluées à celles
réévaluées sur la base de leurs origines et de leurs types d'évolutions
respectives. %
Cette procédure d'appariement est, elle aussi, configurable par le moyen d'une
stratégie de qualité d'appariement. %
Nous offrons ainsi la possibilité de rechercher un appariement exact ou moins
strict dans la règle alternative. %


\section*{Perspectives}
\renewcommand\sectionname{Perspectives}
\markboth{Conclusion}{\sectionname}


Certains des travaux présentés dans ce manuscrit étant en cours d'implantation,
nous disposons déjà d'un certain nombre de perspectives à court terme. %

Tout d'abord, nous prévoyons de terminer l'implantation des déplacements
d'opérations qui peuvent être faites dans une spécification paramétrique et
ainsi compléter la liste de stratégies que nous proposons dans notre modeleur. %
Ainsi, nous pourrons proposer à un utilisateur de changer l'ordre des opérations
dans une spécification
paramétrique. %, mais aussi en supprimer. \david{que le changement d'ordre ne soit pas encore implémenté, ce n'est pas trop grave. Par contre il faudrait que ce soit le cas pour la suppression. Ca te semble jouable d'intégrer ça d'ici la soutenance ?} %
% \xavier{[Remarque : ATTENTION, dans la partie précédente, tu as écrit que le déplacement d'opération est déjà pris en compte (comme l'ajout et la suppression). Mais ici, tu signales que c'est une perspective => contradiction]}
Nous compléterons, dans le même temps, les stratégies de réévaluation lors de
l'ajout d'une suppression d'orbite, de sorte à gérer l'interruption d'un
appariement ou encore la fusion de deux orbites. %

Ensuite, nous poursuivrons l'extension de notre mécanisme de réévaluation pour
permettre la réévaluation de spécifications paramétriques comportant des scripts
Jerboa. %\agnes{pb dans cette phrase}\maxime{mieux ?}%
Nous implanterons les extensions identifiées dans le dernier chapitre pour les
spécifications paramétriques ainsi que les DAG. %
Nous mettrons en \oe uvre des tests afin de contrôler le bon fonctionnement des
procédures de réévaluation. %
Puis, nous étendrons encore nos scripts avec des fonctionnalités telles que des
boucles itérant sur autre chose que des sous-orbites. %

À moyen terme, nous prévoyons d'étendre l'interface graphique de notre modeleur
pour y intégrer des stratégies de réévaluation. %
L'objectif est de permettre aux utilisateurs de configurer leurs modeleurs avec
différentes stratégies, d'en choisir une par défaut ou d'interagir dynamiquement
avec le modeleur au cours de la réévaluation. %
Cette extension concerne les stratégies présentées dans ce manuscrit. %

Pendant la réalisation de ces travaux, nous avons identifié plusieurs
perspectives à plus long terme. %
Celles-ci nous permettent d'ouvrir nos travaux à d'autres champs d'études.

La première piste concerne le domaine de l'archéologie. %
L'objectif est de développer des opérations complexes à l'aide de scripts pour
la modélisation en archéologie. %
En effet, dans le cadre de leur recherche, les archéologues peuvent être amenés
à créer des restitutions virtuelles d'artefacts archéologiques ou des bâtiments
historiques, issus de fouilles et dans des états souvent dégradés. %
Afin de mener à bien ces travaux, ils émettent des hypothèses sur l'état initial
de leur sujet d'étude et ont alors besoin de confronter les différentes
hypothèses de restitution. %
Cette démarche implique ainsi la création d'une variation pour chaque
hypothèse. %
Il s'agit, autrement dit, d'un problème qui peut être pris en compte par la
réévaluation. %
Ainsi, nous prévoyons d'étudier l'applicabilité de notre modeleur et de son
mécanisme de réévaluation pour l'archéologie. %


Une seconde piste que nous avons identifiée, dans la continuité de la première,
est l'annotation temporelle des opérations. %
L'objectif, dans un premier temps, est d'ajouter une annotation temporelle aux opérations. %
Puis, dans un second temps, de permettre la réévaluation du modèle suivant
l'ordre chronologique, c'est-à-dire réappliquer les opérations de la
spécification en fonction des dates qui leur ont été attribuées. %
Pour cela, nous pouvons utiliser les DAG conçus au cours de nos travaux. %
Puisque les DAG représentent des évolutions d'orbites, il est donc possible de les
utiliser pour déterminer l'instant de création d'une orbite et à quel moment
celle-ci est supprimée. %
Ainsi, il est possible de garantir la cohérence des annotations temporelles
vis-à-vis des paramètres topologiques d'une application. %
À terme, une telle extension permettrait de réévaluer des bâtiments historiques, par
exemple, suivant un ordre chronologique supposé. %
Les résultats de ces réévaluations permettant, \textit{in fine}, de confronter
des hypothèses sur la construction et l'évolution de certains biens patrimoniaux
au cours du temps. %
% \david{peut être peux tu ajouter un point avant d'en arriver aux annotations temporelles (histoire d'avoir au moins 2 perspectives à long terme). Je te propose d'ajouter l'idée de développer des opérations complexes à l'aide des scripts dans le domaine de l'archéologie qui a justement un fort besoin de reévaluation pour générer des variantes de ce qui aurait pu exister, car les fouilles ne révèlent que quelques vestiges et il y a une part importante d'hypothèses. Ca nous permet de tester nos scripts et leur réévaluation à un domaine d'application. Si c'est pas très clair, n'hésites pas à m'appeler.}

% Une autre piste envisagée est l'intégration, dans le modeleur, d'annotations à
% partir de données historiques hétérogènes (photographies, notes, mesures,
% textures, propriétés physiques, etc.). %
% Ces annotations seront définies avec un modèle de référence comme
% CIDO-CRM~\cite{cidoccrm}. %
% Une telle contribution nécessiterait alors de garantir la cohérence entre les
% modèles et les annotations. %
% \maxime{Je suis complètement perdu pour celle-ci, on n'a jamais eu le temps d'en
%   parler donc peut-être que je devrais l'enlever ?}

%%% Local Variables:
%%% mode: LaTeX
%%% TeX-master: "../../main"
%%% End:
